% Instrucciones — Reporte Técnico del Proyecto Final
% Geometría Computacional — Universidad Panamericana
% Compilar: pdflatex instrucciones_reporte_tecnico.tex
\documentclass[11pt]{article}

\usepackage[T1]{fontenc}
\usepackage[utf8]{inputenc}
\usepackage[spanish, es-noshorthands]{babel}
\usepackage[a4paper, top=2.5cm, bottom=2.5cm, left=2.8cm, right=2.8cm]{geometry}
\usepackage{enumitem}
\usepackage{parskip}
\usepackage{microtype}
\usepackage{lmodern}
\usepackage{titlesec}
\usepackage[hidelinks]{hyperref}

\setlist{itemsep=1pt, topsep=3pt, parsep=0pt}
\setlist[enumerate]{leftmargin=1.5em}
\setlist[itemize]{leftmargin=1.5em}

\titleformat{\section}[block]{\normalfont\normalsize\bfseries}{}{0pt}{}
\titlespacing{\section}{0pt}{1.8ex plus .3ex}{0.5ex plus .2ex}

\pagestyle{plain}

\begin{document}

\begin{center}
  \textbf{Instrucciones — Reporte Técnico del Proyecto Final}\\[2pt]
  Geometría Computacional \textbullet\ Universidad Panamericana \textbullet\ 2026\\
  {\small Dr.\ Miguel Alcaraz Rivera \textbullet\ malcaraz@up.edu.mx}
\end{center}

\medskip
\hrule
\medskip

\textbf{Objetivo:} Documenta formalmente el sistema implementado en forma de reporte
técnico. Reflexiona y agrega cómo es que tu proyecto aporta a los resultados de
aprendizaje de tu carrera. Agrega esta información a la introducción y conclusiones.

\textbf{Resultados de aprendizaje:}
\begin{itemize}
  \item Diseñar soluciones creativas para problemas complejos de ingeniería.
  \item Diseñar sistemas considerando salud y seguridad públicas, costo total de
        vida, carbono neto cero, y aspectos de recursos, culturales, sociales y
        ambientales.
  \item Crear, seleccionar, aplicar y reconocer las limitaciones de técnicas y
        herramientas modernas de TI, incluyendo matemáticas conceptuales, análisis
        numérico y aspectos formales de informática y ciencias de la información.
\end{itemize}

\medskip
\hrule
\medskip

\section*{Estructura}

El reporte contendrá las siguientes secciones en el orden indicado:

\begin{enumerate}
  \item \textbf{Portada.} Nombre del proyecto, integrantes, materia, institución
        y fecha.
  \item \textbf{Resumen.} Máximo 250 palabras: problema abordado, algoritmos
        empleados y resultados obtenidos.
  \item \textbf{Introducción.} Contexto del problema en GIS y cartografía digital,
        definición precisa del reto de ingeniería, justificación de los algoritmos
        seleccionados, alcance del sistema y organización del documento.
  \item \textbf{Marco teórico.} Investigación bibliográfica sobre los fundamentos
        del sistema (véase la sección correspondiente más adelante).
  \item \textbf{Diseño del sistema.} Arquitectura, módulos, flujo de datos y
        decisiones de diseño con su justificación.  Se incluirá al menos un
        diagrama de componentes o de flujo.
  \item \textbf{Implementación.} Fragmentos de código representativos con
        explicación.  No se incluirá la totalidad del código fuente; se
        seleccionarán los segmentos que ilustren las decisiones de diseño.
  \item \textbf{Resultados y demostración.} Capturas de pantalla requeridas con
        pie de figura.  Descripción de aplicaciones reales del sistema.
  \item \textbf{Conclusiones.} Reflexión sobre las decisiones de diseño, las
        limitaciones encontradas, las implicaciones de aplicación real y el
        trabajo futuro.
  \item \textbf{Bibliografía.} Fuentes académicas en formato IEEE o APA, aplicado
        de manera consistente.  Las fuentes deberán estar citadas en el texto;
        toda afirmación técnica requiere su referencia.
\end{enumerate}

\section*{Marco teórico}

Se cubrirán los siguientes temas obligatorios:

\begin{itemize}
  \item Subdivisiones planares y DCEL: componentes (vértices, semi-aristas,
        caras) y operaciones soportadas.
  \item Paradigma de barrido de plano: línea de barrido, cola de eventos y
        árbol de estado.
  \item Algoritmo de Bentley-Ottmann: tipos de eventos, complejidad en tiempo
        y espacio.
  \item Identificación de caras en una subdivisión planar; cara no acotada.
\end{itemize}

Fuentes recomendadas: de Berg et al., \textit{Computational Geometry}
(3.ª~ed., Springer, 2008), capítulos~2 y~5; Bentley \& Ottmann, \textit{IEEE
Transactions on Computers}, 28(9), 1979.


\section*{Criterios de evaluación}

\begin{itemize}
  \item \textbf{Claridad conceptual:} el sistema puede comprenderse sin consultar
        el código fuente.
  \item \textbf{Marco teórico:} los temas se desarrollan con profundidad y las
        afirmaciones técnicas están respaldadas por referencias.
  \item \textbf{Resultados de aprendizaje:} los tres resultados se abordan de
        forma explícita y argumentada en introducción y conclusiones.
  \item \textbf{Imágenes:} constituyen evidencia del funcionamiento con pie de
        figura explicativo.
  \item \textbf{Honestidad técnica:} las limitaciones del sistema y de los
        algoritmos se reconocen y describen con precisión.
  \item \textbf{Consistencia formal:} notación uniforme, figuras numeradas y
        referenciadas, bibliografía completa.
\end{itemize}

\end{document}
